\documentclass[12pt]{article}
\usepackage[utf8]{inputenc}
\usepackage[T1]{fontenc}
\usepackage{hyperref}
\title{Mattokki Language Resources for K-12 Curriculum Development}
\author{Compiled by Claude Sonnet 5.5 (claude-sonnet-5-5) for Andr\'as Kornai}
\date{}
\begin{document}
\maketitle

\section{Introduction}
Mattokki (ISO 639-3 xnz) is also called Kenzi, Kenuzi or Kunuz. It is a Nubian language of southern Egypt (Aswan area), closely related to Dongolawi/Andaandi of Sudan. Reported speaker numbers vary widely (about 35,000 on Wikipedia, 64,000 in Joshua Project data). Egyptian Arabic and Modern Standard Arabic dominate education and official life. No school programme teaching Mattokki and no dedicated textbook was found in the searches. The results were mostly general reference pages.

\section{Primary Educational Resources}
No primers, textbooks or school materials were found in search results.

\section{Linguistic and Reference Materials}
\begin{itemize}
\item Wikipedia, ``Kenzi language'': \url{https://en.wikipedia.org/wiki/Kenzi_language}
\item Journal article on the linguistic landscape of Mattokki in Aswan, which finds it visible in tourist localities but absent from official signage: \url{https://ttaip.journals.ekb.eg/article_400352.html}
\end{itemize}

\section{Religious and Community Resources}
\begin{itemize}
\item Joshua Project profile, Kunuz Nubian in Egypt (a mission-oriented demographic profile, not a community-run resource): \url{https://joshuaproject.net/people_groups/20525/EG}
\end{itemize}
No scripture translation was confirmed.

\section{Digital and Audio Resources}
No audio, video or software resources were found in search results.

\section{Curriculum Development Context}
The language is spoken in the community, but sources show no presence in schooling. Any curriculum would need to start with new material development together with Nubian community organisations. Orthography (Latin, Arabic and Coptic-derived scripts are all mentioned in sources) would have to be agreed with the community first. Further searching, for example for Nubian cultural centres in Aswan, linguistic documentation of Nubian varieties, and Nubian-language teaching initiatives, is recommended.
\end{document}